% Corte mecánico íntegro: parte_iv.tex:220--336.
% Residencia material del ensamblaje sucesor de 102 capítulos.
% Residencia de realización: capítulo 58.
\chapter{Central electron and helical reader of matter}

\section{Uniqueness of the Material Center under Cellular Inversion}

The material atlas has domain \(\mathcal C_{81}=\{1,\ldots,9\}^2\) and
involution
\[
 \rho_c(r,c)=(10-r,10-c).
\]
Its only fixed point is
\[
 (r,c)=(5,5).
\]
The other \(80\) cells form \(40\) pairs. The central electron is the
only point multisection; the remaining internal species are fibers, orbits,
or multisections. Spiral theory does not call this conclusion into question.

\section{Common tritic mode}

In the canonical APP chart \(1,\ldots,9\), the vector of values of the operational
tritic phase selector is
\[
 m_{\ph}=(1,-1,0)^3.
\]
For the three-dimensional joint table,
\[
 C^{(3)}m_{\ph}=(C^{(3)})^{\!*}m_{\ph}
 =-\frac12m_{\ph}.
\]
Therefore
\[
 s=\frac12,\qquad
 1-s^2=\frac34=\frac{90}{120},
\]
\begin{equation}
 \|[P_{\Sigma,3},P_{\Pi,3}]\|
 =\frac{\sqrt3}{4}.
 \label{espiral:eq:modo-electron-app}
\end{equation}
The mode, oriented defect \(90\to120\), and electronic prefactor are
coordinated publications of the same APP--TRIT--TPK chain, whose effective
composition provides their internal intertwining.

\section{Ledger material}

Each cell of the atlas bears the signature
\[
 \mathcal R(r,c)=
 (n_A,s_C,k_{\Delta_4},\nu_{120},\nu_{270},\tau_{12})_{r,c}.
\]
The image contains \(56\) route families; refinement by family and
level contains \(13\) multisections. The governing counts are
\[
 81/56/13/324,
\]
with the null sector and the global towers. The radial reader only can be added
by means of a compatible signature, by ejemplo
\[
 \Sigma_{\mathrm{esp}}(X)=
 (p_X,a_X,[\Lambda_X,C_X],
 [\kappa_{9,X}],\mu_X,\varepsilon_{{\mathrm{tw}},X}),
\]
that is constant where appropriate and natural under refinement.

\section{Material compatibility square}

Let \(\mathcal E_{\mathrm{atlas}}\) the space of persistent sections and
\(\mathfrak S_{\mathrm{esp}}\) the spiral reader. The material compatibility
is expressed by means of the square
\[
\begin{array}{ccc}
\mathcal E_{\mathrm{atlas}}
 & \xrightarrow{\ \mathfrak S_{\mathrm{esp}}\ }
 & \mathcal D_{\rad}^{\enr}\\[1.0em]
{\scriptstyle (K_D,K_\Omega)}\big\downarrow
 && \big\downarrow{\scriptstyle \ev_{\mathrm{dim}}^{\mathrm{esp}}}\\[1.0em]
\operatorname{Spec}(K_D,K_\Omega)
 & \xrightarrow{\ \mathfrak R_{\mathrm{mat}}\ }
 & \mathcal M_{\mathrm{hel}}
\end{array}
\]
The electron \((5,5)\) provides the central anchor, but does not select
\(p,\Lambda,C\) retrospectively. The other species require
constancy on multisections, compatibility with cell inversion,
\(K_D/K_\Omega\), CKM/PMNS, color, and spin representation.

\section{Spin and helicity}

is not identified the spin \(1/2\) with a vuelta of a helix. The
fermionic representation satisfies
\[
 U(2\pi)\psi_e=-\psi_e,
 \qquad
 U(4\pi)\psi_e=\psi_e.
\]
The nonadic helix records phase and memory; the spin representation lives
in another codomain. A valid material intertwiner must preserve both
laws, not replace one with the other.

\begin{figure}[htbp]
\centering
\begin{tikzpicture}[scale=.63]
\foreach \i in {1,...,9}{
 \foreach \j in {1,...,9}{
  \draw[HMTLine,fill=HMTLight] (\i,\j) rectangle +(1,1);
 }
}
\fill[HMTNavy] (5.15,5.15) rectangle (5.85,5.85);
\draw[HMTRed,very thick] (5.5,5.5) circle (1.05);
\draw[->,HMTRed,thick] (5.5,5.5)--(8.9,8.9);
\node[right,HMTNavy,font=\sffamily\bfseries] at (8.9,8.9) {unique fixed point \((5,5)\)};
\node[below,HMTGray,font=\small\sffamily] at (5.5,.45) {atlas material \(9\times9\)};
\end{tikzpicture}
\caption{The central electron is a result of the atlas; the spiral correspondence is a subsequent reader.}
\label{espiral:fig:atlas-electron}
\end{figure}
